\section{Managing Team Issues}
\textit{Reference $|$ Course Text: Chapter 13}

Team dysfunction is the most common source of PA intervention requests. Most issues are predictable and manageable if identified early. This section provides frameworks for diagnosis, escalation, and resolution.

\subsection{The Tuckman Model Applied to Capstone}

Teams progress through Forming, Storming, Norming, and Performing. In Capstone Design, the timeline is compressed:

\begin{itemize}[nosep,leftmargin=1.5em]
\item \textbf{Forming (Sprints 0--1):} Polite, cautious, low conflict. Teams avoid disagreement to maintain harmony. This is normal.
\item \textbf{Storming (Sprints 3--5):} Conflict emerges as the workload increases and design decisions require trade-offs. Personality differences surface. This is the highest-risk phase.
\item \textbf{Norming (Sprints 5--8):} Team establishes working norms. Roles solidify. Conflict is managed, not eliminated. PDR often serves as a unifying forcing function.
\item \textbf{Performing (Sprints 9--14):} Effective teams reach this phase and execute efficiently. Struggling teams may never leave Storming.
\end{itemize}

\begin{bestpractice}
Storming is healthy and expected. Your role during Storming is not to eliminate conflict but to ensure it stays productive. Productive conflict is about design decisions and priorities. Destructive conflict is personal, disrespectful, or passive-aggressive. Intervene in the latter; facilitate the former.
\end{bestpractice}

\subsection{Identifying Underperformers Early}

Watch for these indicators starting in Sprint~2:

\begin{itemize}[nosep,leftmargin=1.5em]
\item \textbf{Missed deadlines.} One slip is human. A pattern of two or more missed sprint deadlines signals disengagement.
\item \textbf{Absent from meetings.} Missing team meetings or client meetings without advance notice.
\item \textbf{No commits or contributions.} Version history and repository logs show minimal or zero contributions from the individual.
\item \textbf{Cannot explain their own work.} When asked about their assigned subsystem at a sprint review, they defer to another team member.
\item \textbf{Peers stop relying on them.} Other team members quietly redistribute the underperformer's work without telling you. This often surfaces in peer evaluations.
\end{itemize}

\begin{redflag}
If a team member has contributed nothing visible in two consecutive sprints, do not wait for the next peer evaluation. Address it directly with the individual in a one-on-one conversation.
\end{redflag}

\subsection{Escalation Path}

Follow this escalation sequence. Each step should be documented.

\begin{enumerate}[nosep,leftmargin=1.5em]
\item \textbf{Team conversation.} Encourage the team to address the issue internally first. Healthy teams can resolve most issues through direct conversation. Ask the team: ``Have you talked to [Name] about this directly?''

\item \textbf{PA-mediated conversation.} If the team conversation fails or the team is unwilling to have it, facilitate a structured conversation. Meet with the team (or the individuals involved) and establish specific, measurable expectations for the next sprint. Document the agreement.

\item \textbf{PA one-on-one.} For persistent issues, meet privately with the underperforming student. Be direct: ``Your teammates and I have observed [specific behaviors]. Here is what needs to change by [specific date].'' Document this conversation.

\item \textbf{CF intervention.} If the PA-mediated conversation does not produce change within one sprint, escalate to the CF. The CF may initiate a Performance Improvement Plan (PIP), reassign team roles, or take other corrective action.
\end{enumerate}

\begin{paaction}
Document every escalation step with dates, attendees, and specific commitments made. This documentation protects you, the student, and the program if the situation becomes a formal dispute.
\end{paaction}

\subsection{The Bus Factor Problem}

The ``bus factor'' is the number of team members who could be unavailable before the project stalls. A bus factor of one means all critical knowledge resides in a single person.

\textbf{Symptoms:}
\begin{itemize}[nosep,leftmargin=1.5em]
\item Only one person can explain the firmware, the CAD model, or the test procedure.
\item One team member is CC'd on every client email; others are not.
\item Team meetings cannot proceed when a specific member is absent.
\end{itemize}

\textbf{Remediation:}
\begin{itemize}[nosep,leftmargin=1.5em]
\item Require cross-training: each subsystem should have a primary and a secondary owner.
\item Require documentation of all critical procedures (not just the final report---working documents).
\item At sprint reviews, ask questions to the secondary owner, not the primary, to test knowledge transfer.
\end{itemize}

\subsection{When to Involve the CLT Director}

Escalate to the CLT Director (not just the CF) when:

\begin{itemize}[nosep,leftmargin=1.5em]
\item The issue involves the client relationship (client is unresponsive, client is making unreasonable demands, client contact has changed).
\item The issue involves safety concerns in fabrication spaces (InnoHub, labs).
\item The issue involves facilities, equipment access, or procurement logistics beyond normal channels.
\item A team member raises a concern about harassment, discrimination, or hostile behavior (the CLT Director can involve institutional resources).
\end{itemize}

\subsection{Personality Conflicts vs.\ Performance Issues}

Distinguish between these two categories, as the intervention differs:

\begin{faqbox}[Personality conflict vs.\ performance issue]
\textbf{Personality conflict:} All members are contributing, but interpersonal friction makes collaboration difficult. Two members disagree on approach and the disagreement has become personal.\\
\textit{Intervention:} Facilitate a structured conversation focused on shared goals. Establish communication norms. Remind the team that they do not need to be friends; they need to be professional.

\textbf{Performance issue:} One or more members are not meeting commitments, regardless of interpersonal dynamics. Deadlines are missed, work quality is below standard, or attendance is poor.\\
\textit{Intervention:} Follow the escalation path above. Performance issues require specific, measurable expectations and consequences.
\end{faqbox}

\begin{tipbox}
When a team reports a ``personality conflict,'' investigate whether it is actually a performance issue in disguise. Often, interpersonal friction originates from one member not carrying their share of the workload, and the resulting resentment manifests as personal conflict.
\end{tipbox}


% ═══════════════════════════════════════
